% =============================================================================
% Rozdział 3: Projekt architektury systemu badawczego
% =============================================================================
\chapter{Projekt architektury systemu badawczego}

Rozdział niniejszy opisuje projekt architektury systemu badawczego — jego komponenty, strategie skalowania, infrastrukturę monitorującą oraz plan eksperymentów. Wszystkie decyzje projektowe zostały podporządkowane celowi badawczemu: porównaniu skuteczności trzech strategii \textit{HorizontalPodAutoscaler} (\textit{HPA}) dla dwóch diametralnie różnych profili obciążeniowych.

\section{Macierz badawcza}

Podstawą metodologii badawczej jest macierz $2 \times 3$ obejmująca dwie aplikacje o odmiennych profilach obciążeniowych i trzy strategie HPA:

\begin{table}[H]
  \centering
  \begin{tabularx}{\textwidth}{|l|X|X|X|}
    \hline
    \textbf{Profil aplikacji} & \textbf{HPA CPU (50\%)} & \textbf{HPA Memory (70\%)} & \textbf{HPA Custom RPS (100/Pod)} \\
    \hline
    \textit{CPU-bound} (cpu-service) & \texttt{cpu-cpu} & \texttt{cpu-memory} & \texttt{cpu-custom} \\
    \hline
    \textit{I/O-bound} (io-service) & \texttt{io-cpu} & \texttt{io-memory} & \texttt{io-custom} \\
    \hline
  \end{tabularx}
  \caption{Macierz badawcza $2 \times 3$ — identyfikatory scenariuszy testowych}
  \label{tab:research-matrix}
\end{table}

Każda z sześciu kombinacji jest testowana w odrębnym scenariuszu, co daje łącznie sześć scenariuszy właściwych (plus dwa scenariusze bazowe bez HPA). Macierz umożliwia izolowanie wpływu pojedynczej zmiennej niezależnej (strategia HPA) przy kontrolowanym profilu obciążeniowym.

Oprócz sześciu scenariuszy macierzy badawczej zdefiniowano dwa scenariusze bazowe (\textit{baseline}):

\begin{itemize}
  \item \texttt{baseline-cpu} — \textit{CPU-bound}, brak HPA, 1 replika na stałe.
  \item \texttt{baseline-io} — \textit{I/O-bound}, brak HPA, 1 replika na stałe.
\end{itemize}

Scenariusze bazowe służą jako punkt odniesienia do oceny, czy HPA w ogóle poprawia wydajność w stosunku do statycznej konfiguracji. Dodatkowo, dla testów integracyjnych zdefiniowano scenariusz łączony (\texttt{combined}), w którym oba serwisy są obciążane jednocześnie.

\section{Aplikacje badawcze}

System składa się z trzech mikroserwisów — dwóch podlegających badaniu (\texttt{cpu-service}, \texttt{io-service}) oraz jednego pomocniczego (\texttt{echo-server}). Każdy mikroserwis jest niezależną aplikacją \textit{FastAPI} (Python 3.12) konteneryzowaną za pomocą \textit{Docker} i wdrażaną jako \textit{Deployment} w \textit{Kubernetes}.

\subsection{CPU-service — obciążenie CPU-bound}

\texttt{cpu-service} implementuje dwa endpointy obciążające procesor:

\begin{itemize}
  \item \texttt{POST /process} — przetwarzanie obrazu: akceptuje plik graficzny (\textit{multipart form}), wykonuje sekwencję operacji CPU-intensywnych z użyciem biblioteki \textit{Pillow} (m.in. konwersja do skali szarości, filtr \texttt{BLUR}, filtr \texttt{SHARPEN}, filtr \texttt{FIND\_EDGES}, ekstrakcja kanałów RGB, resize), oblicza entropię Shannona dla każdego kanału, zwraca statystyki przetwarzania (wymiary, entropia, czas).
  \item \texttt{GET /fibonacci} — obliczanie $n$-tej liczby Fibonacciego rekurencyjnym algorytmem o złożoności $O(2^{n})$. Parametr $n$ jest konfigurowalny (domyślnie $n = 25$).
\end{itemize}

Obie operacje są czysto CPU-intensywne — nie wykonują operacji sieciowych ani dyskowych poza odbiorem żądania i wysłaniem odpowiedzi. Dzięki temu zużycie CPU silnie koreluje z obciążeniem żądaniami.

\subsection{IO-service — obciążenie I/O-bound}

\texttt{io-service} implementuje dwa endpointy symulujące operacje wejścia-wyjścia:

\begin{itemize}
  \item \texttt{GET /query} — łańcuch asynchronicznych opóźnień: wykonuje sekwencję wywołań \texttt{asyncio.sleep()} z konfigurowalnym opóźnieniem bazowym (\texttt{delay}, domyślnie 100\,ms), liczbą kroków (\texttt{steps}, domyślnie 3) oraz losowym rozrzutem (\textit{jitter}, domyślnie $\pm20$\%). Całkowity czas odpowiedzi jest sumą opóźnień, a procesor pozostaje bezczynny podczas oczekiwania.
  \item \texttt{GET /upstream} — wywołanie zewnętrzne: przekazuje żądanie do \texttt{echo-server} i zwraca odpowiedź z dodatkowym opóźnieniem symulującym przetwarzanie.
\end{itemize}

Serwis wykorzystuje \texttt{asyncio.Semaphore(50)} do ograniczenia liczby współbieżnych zapytań, co zapobiega wyczerpaniu zasobów systemowych przy bardzo wysokim obciążeniu.

\subsection{Echo-server — serwis pomocniczy}

\texttt{echo-server} jest minimalistycznym serwisem pełniącym rolę zależności zewnętrznej dla \texttt{io-service}. Udostępnia dwa endpointy:

\begin{itemize}
  \item \texttt{GET /health} — \textit{health check}.
  \item \texttt{POST /echo} — zwraca kopię otrzymanego żądania (nagłówki, body, parametry), umożliwiając weryfikację poprawnego przekazywania danych.
\end{itemize}

\textit{Echo-server} jest wyłączony z macierzy badawczej — służy wyłącznie jako cel wywołań w scenariuszu \texttt{external\_calls} dla \texttt{io-service}.

\subsection{Memory-service — serwis pamięciowy (dodatkowy)}

W ramach rozwoju projektu zaimplementowano również \texttt{memory-service} — serwis z pamięciową bazą klucz-wartość (\textit{in-memory cache}) symulujący obciążenie typu \textit{memory-bound}. Udostępnia endpointy \texttt{GET/POST/DELETE /cache}, \texttt{POST /cache/fill} (wypełnienie cache danymi testowymi o kontrolowanym rozmiarze) oraz \texttt{GET /stats} (statystyki trafień). \texttt{memory-service} nie został włączony do głównej macierzy badawczej na obecnym etapie prac — stanowi rozszerzenie do potencjalnych przyszłych badań.

\section{Strategie HPA}

Dla każdej aplikacji badawczej zdefiniowano trzy warianty \textit{HorizontalPodAutoscaler}, różniące się wyłącznie typem metryki. Wszystkie HPA współdzielą następujące parametry:

\begin{itemize}
  \item \texttt{minReplicas}: 2
  \item \texttt{maxReplicas}: 10
  \item \texttt{apiVersion}: \texttt{autoscaling/v2}
\end{itemize}

\subsection{HPA CPU (Resource Metric)}

Strategia oparta na zużyciu CPU — HPA monitoruje średnie zużycie procesora jako procent \textit{resource request} dla wszystkich Podów w \textit{Deploymencie}. Wartość progowa: 50\%. Oznacza to, że jeśli średnie zużycie CPU przekracza 50\% \texttt{requests} (czyli 125\,m dla cpu-service, 50\,m dla io-service), HPA dodaje repliki.

\lstinputlisting[
  language=yaml,
  caption={Konfiguracja HPA CPU dla \texttt{cpu-service} (plik \texttt{k8s/cpu-service/hpa-cpu.yaml})},
  label={lst:hpa-cpu-config},
  firstline=1,
  lastline=22
]{../k8s/cpu-service/hpa-cpu.yaml}

Strategia CPU jest najczęściej stosowaną i domyślnie rekomendowaną w ekosystemie \textit{Kubernetes}. Jej skuteczność zależy od tego, czy zużycie CPU koreluje z obciążeniem aplikacji.

\subsection{HPA Memory (Resource Metric)}

Strategia oparta na zużyciu pamięci — HPA monitoruje średnie zużycie pamięci jako procent \textit{resource request}. Wartość progowa: 70\%. Strategia ta jest rzadziej stosowana, ponieważ większość aplikacji alokuje pamięć w sposób nieelastyczny (po uruchomieniu lub przy pierwszym żądaniu), co sprawia, że metryka pamięci nie odzwierciedla bieżącego obciążenia.

\lstinputlisting[
  language=yaml,
  caption={Konfiguracja HPA Memory dla \texttt{cpu-service} (plik \texttt{k8s/cpu-service/hpa-memory.yaml})},
  label={lst:hpa-memory-config},
  firstline=1,
  lastline=23
]{../k8s/cpu-service/hpa-memory.yaml}

Hipoteza badawcza przewiduje, że HPA Memory będzie nieskuteczne dla obu profili obciążeniowych — ani CPU-bound, ani I/O-bound nie powinny generować istotnych wahań zużycia pamięci skorelowanych z ruchem.

\subsection{HPA Custom RPS (Pods Metric)}

Strategia oparta na niestandardowej metryce \texttt{http\_requests\_per\_second} — liczbie żądań HTTP na sekundę na Pod. Metryka jest dostarczana przez \textit{prometheus-adapter}, który transformuje liczniki \textit{Prometheusa} (\texttt{cpu\_service\_http\_requests\_total}, \texttt{io\_service\_http\_requests\_total}) na chwilową wartość RPS za pomocą funkcji \texttt{rate()\lbrack{}1m\rbrack{}}. Wartość progowa: 100 żądań na sekundę na Pod.

\lstinputlisting[
  language=yaml,
  caption={Konfiguracja HPA Custom RPS (plik \texttt{k8s/cpu-service/hpa-custom.yaml})},
  label={lst:hpa-custom-config},
  firstline=1,
  lastline=24
]{../k8s/cpu-service/hpa-custom.yaml}

Strategia RPS jest bezpośrednio powiązana z obciążeniem żądaniami — niezależnie od tego, czy żądania są CPU-intensywne, czy I/O-intensywne, liczba żądań na sekundę rośnie proporcjonalnie do ruchu. Dzięki temu HPA Custom RPS powinna być skuteczna dla obu profili obciążeniowych.

\section{Architektura wdrożeniowa w GCP}

\subsection{Topologia sieciowa GCP}

System badawczy jest w całości wdrożony w środowisku \textit{Google Cloud Platform} (GCP). Wszystkie komponenty znajdują się w jednej sieci \textit{Virtual Private Cloud} (VPC), co zapewnia niskie opóźnienia i bezpieczną komunikację wewnętrzną:

\begin{figure}[H]
  \centering
  % Opis architektury w formie diagramu tekstowego (TikZ może być dodany później)
  \begin{verbatim}
  +--------------------------------------------------+
  |                GCP VPC Network                    |
  |                                                   |
  |  +------------------+    +---------------------+  |
  |  |  k6-generator    |    |   GKE Cluster       |  |
  |  |  (Compute Engine)|    |                     |  |
  |  |                  |    |  +---------------+  |  |
  |  |  k6 -> LB IP     |--->|  | LoadBalancer  |  |  |
  |  |  (hosts mapping) |    |  +-------+-------+  |  |
  |  |                  |    |          |          |  |
  |  +------------------+    |  +-------v-------+  |  |
  |                          |  | cpu-service   |  |  |
  |                          |  | io-service    |  |  |
  |                          |  | echo-server   |  |  |
  |                          |  +---------------+  |  |
  |                          |                     |  |
  |                          |  kube-prometheus-   |  |
  |                          |  stack (Helm)       |  |
  |                          |  +---------------+  |  |
  |                          |  | Prometheus    |  |  |
  |                          |  | Grafana (LB)  |  |  |
  |                          |  +---------------+  |  |
  |                          +---------------------+  |
  +--------------------------------------------------+
  \end{verbatim}
  \caption{Topologia systemu badawczego w GCP — wszystkie komponenty w jednej sieci VPC}
  \label{fig:gcp-topology}
\end{figure}

\subsection{GKE Cluster — klaster Kubernetes}

Klaster GKE jest centrum systemu badawczego. Wszystkie aplikacje, HPA, monitoring i \textit{prometheus-adapter} działają jako zasoby \textit{Kubernetes} w namespace \texttt{magisterka}. Klaster składa się z węzłów typu \texttt{e2-standard-4} (4 vCPU, 16\,GB RAM), które zapewniają wystarczające zasoby dla maksymalnie 20 Podów aplikacyjnych (przy \texttt{maxReplicas=10} dla obu serwisów) plus Podów monitorujących.

\subsection{GCP LoadBalancer — ekspozycja serwisów}

Każdy serwis aplikacyjny (\texttt{cpu-service}, \texttt{io-service}, \texttt{echo-server}) jest eksponowany jako \texttt{Service} typu \texttt{LoadBalancer}. GKE automatycznie tworzy zewnętrzny \textit{GCP LoadBalancer} (Layer 4, TCP) i przydziela mu publiczny adres IP. Ruch z maszyny \texttt{k6-generator} trafia bezpośrednio na te adresy IP.

Dodatkowo, \textit{Grafana} (część \textit{kube-prometheus-stack}) jest również eksponowana przez \texttt{LoadBalancer}, co umożliwia obserwację eksperymentów w czasie rzeczywistym z dowolnej przeglądarki.

\subsection{k6-generator — maszyna generująca obciążenie}

Dedykowana maszyna \textit{Compute Engine} (\texttt{k6-generator}) w tej samej sieci VPC uruchamia skrypty \textit{k6}. Mapowanie nazw serwisów na adresy IP \textit{LoadBalancer} odbywa się przez mechanizm \texttt{options.hosts} w skryptach \textit{k6} — bez konieczności konfiguracji DNS:

\begin{lstlisting}[language=javascript, caption={Mapowanie hostów w skrypcie \textit{k6} — zamiast DNS}, label={lst:k6-hosts-arch}]
export const options = {
  hosts: {
    'cpu-service': '34.118.71.122',    // GCP LoadBalancer IP
    'io-service': '35.204.93.45',      // GCP LoadBalancer IP
  },
  scenarios: { ... }
};
\end{lstlisting}

To rozwiązanie eliminuje potrzebę konfiguracji serwera DNS, rekordów \texttt{/etc/hosts} czy tunelowania przez \texttt{kubectl port-forward}. Ruch sieciowy między \texttt{k6-generator} a Podami w GKE przechodzi przez GCP LoadBalancer, który dystrybuuje go do zdrowych Podów.

\subsection{Rejestr obrazów kontenerowych}

Obrazy Docker są przechowywane w \textit{GCP Artifact Registry}, skąd są pobierane przez GKE podczas wdrażania. Eliminuje to konieczność ręcznego importowania obrazów do \textit{container runtime} na każdym węźle — GKE automatycznie pobiera obrazy z rejestru na podstawie pola \texttt{image} w \textit{Deploymencie}.

\section{Infrastruktura monitorująca}

\subsection{kube-prometheus-stack}

Pakiet \textit{kube-prometheus-stack} jest instalowany przez \textit{Helm} i dostarcza kompletny stos monitorujący:

\begin{itemize}
  \item \textbf{Prometheus} — zbiera metryki infrastrukturalne (przez \texttt{kube-state-metrics}, \texttt{node-exporter}, \texttt{cAdvisor}) oraz aplikacyjne (z endpointów \texttt{/metrics} serwisów) co 15 sekund.
  \item \textbf{Grafana} — udostępnia dashboardy wizualizacyjne, eksponowana publicznie przez GCP LoadBalancer dla obserwacji eksperymentów w czasie rzeczywistym.
  \item \textbf{kube-state-metrics} — generuje metryki o stanie obiektów \textit{Kubernetes}: Deploymenty, HPA, Pody.
  \item \textbf{node-exporter} — zbiera metryki systemowe z węzłów (CPU, pamięć, dysk, sieć).
\end{itemize}

\textit{Prometheus Operator} zarządza cyklem życia \textit{Prometheusa} i automatycznie konfiguruje \textit{scraping} na podstawie obiektów \textit{ServiceMonitor} i \textit{PodMonitor}. Aplikacje badawcze są automatycznie wykrywane dzięki adnotacjom \texttt{prometheus.io/scrape: "true"} na Podach.

\subsection{Grafana — dashboard i eksport danych}

\textit{Grafana} udostępnia autorski dashboard ,,Magisterka — System Overview'' zawierający panele wizualizujące kluczowe metryki systemu w czasie rzeczywistym:

\begin{itemize}
  \item RPS dla \texttt{cpu-service} i \texttt{io-service} (panele typu \texttt{stat} z miniwykresem).
  \item Zużycie CPU na Pod (panele typu \texttt{timeseries}).
  \item Zużycie pamięci na Pod.
  \item Liczba replik per \textit{Deployment}.
  \item Opóźnienia (\textit{latency}): p50, p95, p99.
  \item Liczba żądań w toku (\textit{in-flight}).
  \item Stan HPA (pożądana vs. aktualna liczba replik).
\end{itemize}

Dla każdego z 8 scenariuszy testowych eksportowanych jest 5 zestawów danych CSV z \textit{Grafany}:

\begin{enumerate}
  \item \texttt{replicas.csv} — liczba replik HPA w czasie.
  \item \texttt{cpu\_mcpu.csv} — zużycie CPU per Pod w milijednostkach (mCPU).
  \item \texttt{memory\_mb.csv} — zużycie pamięci per Pod w megabajtach.
  \item \texttt{rps.csv} — przepustowość (RPS) w czasie.
  \item \texttt{latency\_ms.csv} — opóźnienia odpowiedzi w czasie (p50, p95, p99).
\end{enumerate}

Eksporty te, wraz z wynikami \textit{k6} (JSON + tekstowe podsumowanie), stanowią kompletny zbiór danych dla każdego scenariusza.

\subsection{Prometheus Adapter}

\textit{Prometheus Adapter} (wersja \texttt{v0.12.0}) implementuje \textit{Custom Metrics API} w \textit{Kubernetes}, umożliwiając HPA odczytywanie metryk z \textit{Prometheusa}. Konfiguracja adaptera (przechowywana w \texttt{ConfigMap} \texttt{prometheus-adapter-config}) definiuje cztery reguły mapowania:

\begin{enumerate}
  \item \texttt{cpu\_service\_http\_requests\_total} $\rightarrow$ \texttt{http\_requests\_per\_second} (przez \texttt{rate()\lbrack{}1m\rbrack{}}).
  \item \texttt{io\_service\_http\_requests\_total} $\rightarrow$ \texttt{http\_requests\_per\_second} (przez \texttt{rate()\lbrack{}1m\rbrack{}}).
  \item \texttt{cpu\_service\_http\_requests\_inprogress} $\rightarrow$ \texttt{http\_requests\_inflight} (wartość chwilowa).
  \item \texttt{io\_service\_http\_requests\_inprogress} $\rightarrow$ \texttt{http\_requests\_inflight} (wartość chwilowa).
\end{enumerate}

Reguły mapowania wykorzystują mechanizm \texttt{<<.LabelMatchers>>} i \texttt{<<.GroupBy>>} do dynamicznego filtrowania po namespace i Podzie, co jest wymagane przez \textit{Custom Metrics API}.

\section{Topologia zasobów Kubernetes}

Wszystkie komponenty systemu są wdrażane w pojedynczym \textit{namespace} \texttt{magisterka}. Topologia obejmuje:

\begin{itemize}
  \item \textbf{3 Deploymenty} aplikacyjne: \texttt{cpu-service}, \texttt{io-service}, \texttt{echo-server} (plus opcjonalnie \texttt{memory-service}).
  \item \textbf{3 Serwisy} typu \texttt{LoadBalancer}: każdy kieruje ruch na port 80 do \texttt{targetPort 8000} odpowiedniego \textit{Deploymentu}, eksponując serwis przez GCP LoadBalancer.
  \item \textbf{6 HPA} (po 3 na serwis badawczy): pliki \texttt{hpa-cpu.yaml}, \texttt{hpa-memory.yaml}, \texttt{hpa-custom.yaml}.
  \item \textbf{Monitoring}: \textit{kube-prometheus-stack} (Prometheus + Grafana + kube-state-metrics + node-exporter) + \texttt{prometheus-adapter}.
\end{itemize}

W środowisku testowym ruch z \textit{k6} do serwisów jest kierowany przez GCP LoadBalancer na podstawie mapowania \texttt{options.hosts}. Opóźnienie sieciowe między maszyną \texttt{k6-generator} a Podami w GKE wynosi $\sim$1--5\,ms (komunikacja wewnątrz tej samej strefy VPC).

\subsection{Zasoby przydzielone serwisom}

\begin{table}[H]
  \centering
  \begin{tabularx}{\textwidth}{|l|X|X|X|}
    \hline
    \textbf{Serwis} & \textbf{CPU request} & \textbf{CPU limit} & \textbf{Memory request / limit} \\
    \hline
    \texttt{cpu-service} & 250\,m & 1000\,m & 256\,Mi / 512\,Mi \\
    \texttt{io-service} & 100\,m & 500\,m & 128\,Mi / 256\,Mi \\
    \texttt{echo-server} & 50\,m & 200\,m & 64\,Mi / 128\,Mi \\
    \texttt{memory-service} & 200\,m & 500\,m & 256\,Mi / 1024\,Mi \\
    \hline
  \end{tabularx}
  \caption{Zasoby przydzielone serwisom (\textit{resource requests} i \textit{limits})}
  \label{tab:resource-allocation}
\end{table}

Różnica w przydziale zasobów między \texttt{cpu-service} a \texttt{io-service} odzwierciedla odmienne profile obciążeniowe: \texttt{cpu-service} potrzebuje więcej CPU (250\,m vs. 100\,m) i pamięci (256\,Mi vs. 128\,Mi) ze względu na operacje przetwarzania obrazu.

\section{Plan eksperymentów}

\subsection{Scenariusze testowe}

Eksperyment obejmuje 8 scenariuszy uruchamianych sekwencyjnie:

\begin{enumerate}
  \item \texttt{baseline-cpu} — CPU-bound, brak HPA, 1 replika.
  \item \texttt{baseline-io} — I/O-bound, brak HPA, 1 replika.
  \item \texttt{cpu-cpu} — CPU-bound + HPA CPU (50\%).
  \item \texttt{cpu-memory} — CPU-bound + HPA Memory (70\%).
  \item \texttt{cpu-custom} — CPU-bound + HPA Custom RPS (100/Pod).
  \item \texttt{io-cpu} — I/O-bound + HPA CPU (50\%).
  \item \texttt{io-memory} — I/O-bound + HPA Memory (70\%).
  \item \texttt{io-custom} — I/O-bound + HPA Custom RPS (100/Pod).
\end{enumerate}

\subsection{Procedura eksperymentalna}

Każdy scenariusz jest wykonywany według ściśle określonej procedury:

\begin{enumerate}
  \item \textbf{Przełączenie HPA} — usunięcie wszystkich istniejących HPA dla docelowego serwisu i nałożenie pliku HPA odpowiadającego testowanej strategii (przez \texttt{kubectl apply}).
  \item \textbf{Reset Deploymentu} — skalowanie \textit{Deploymentu} do 1 repliki i oczekiwanie na gotowość (\texttt{rollout status}).
  \item \textbf{Uruchomienie testu k6} — wykonanie skryptu obciążeniowego z maszyny \texttt{k6-generator}, ruch kierowany na adres IP GCP LoadBalancer przez mapowanie \texttt{options.hosts}.
  \item \textbf{Zapis wyników} — wyniki k6 (JSON + tekstowe podsumowanie) zapisywane do katalogu \texttt{results/<scenariusz>/}, wraz z 5 eksportami CSV z \textit{Grafany}.
  \item \textbf{Powtórzenie} — kroki 1--4 powtarzane 5-krotnie dla każdego scenariusza.
  \item \textbf{Cooldown} — po każdym scenariuszu: skalowanie obu serwisów do 1 repliki, usunięcie wszystkich HPA, odczekanie 120 sekund na stabilizację klastra.
\end{enumerate}

\subsection{Profile obciążeniowe k6}

Każdy scenariusz wykorzystuje dedykowany skrypt \textit{k6} dopasowany do profilu obciążeniowego:

\textbf{CPU-bound} (skrypt \texttt{cpu-load-test.js}):
\begin{itemize}
  \item \textbf{Scenariusz Fibonacci}: ramp-up od 1 do 50 VU w ciągu 60 sekund, następnie stałe obciążenie 50 VU przez 120 sekund.
  \item \textbf{Scenariusz przetwarzanie obrazu}: stałe 10 VU przez 180 sekund, każde żądanie wysyła losowo wygenerowany obraz PNG ($400 \times 300$ pikseli) do \texttt{POST /process}.
\end{itemize}

\textbf{I/O-bound} (skrypt \texttt{io-load-test.js}):
\begin{itemize}
  \item \textbf{Scenariusz stałych zapytań}: ramp-up od 1 do 100 VU w ciągu 120 sekund, stałe obciążenie przez 180 sekund, \texttt{delay=200\,ms}, \texttt{steps=3}.
  \item \textbf{Scenariusz stres-test}: stałe 30 VU, \texttt{delay} rosnący od 50\,ms do 2000\,ms w krokach co 120 sekund.
  \item \textbf{Scenariusz wywołań zewnętrznych}: stałe 20 VU, \texttt{external\_call=true}, żądania do \texttt{GET /query}.
\end{itemize}

\textbf{Scenariusz łączony} (skrypt \texttt{combined-scenario.js}):
\begin{itemize}
  \item Fibonacci: 30 VU (ramp-up 30\,s, steady 180\,s).
  \item Przetwarzanie obrazu: 5 VU (stałe, 180\,s).
  \item Zapytania IO: 80 VU (ramp-up 60\,s, steady 180\,s).
\end{itemize}

Wszystkie skrypty definiują progi akceptacji (\textit{thresholds}): $p(95) < 5000$\,ms dla scenariuszy CPU, $p(95) < 10000$\,ms dla scenariuszy IO, oraz \texttt{http\_req\_failed} $< 1\%$ (CPU) / $< 5\%$ (IO).

\subsection{Metryki badawcze}

Dla każdego scenariusza zbierane są następujące metryki:

\begin{itemize}
  \item \textbf{Opóźnienie (\textit{latency})} — $p(50)$, $p(95)$, $p(99)$ czasu odpowiedzi HTTP (ms).
  \item \textbf{Przepustowość (\textit{throughput})} — liczba żądań na sekundę (RPS) obsłużonych przez system.
  \item \textbf{Czas do skalowania (\textit{time-to-scale})} — czas od przekroczenia progu HPA do osiągnięcia docelowej liczby replik.
  \item \textbf{Efektywność skalowania (\textit{scaling efficiency})} — stosunek zmiany przepustowości do zmiany liczby replik.
  \item \textbf{Stopa błędów (\textit{error rate})} — odsetek żądań zakończonych kodem HTTP $\geq 400$.
  \item \textbf{Liczba replik w czasie} — przebieg czasowy skalowania (dane z \textit{Prometheusa} przez \textit{Grafana} CSV).
  \item \textbf{Zużycie CPU i pamięci} — średnie i szczytowe zużycie zasobów na Pod (dane z \textit{Grafana} CSV).
\end{itemize}

\subsection{Analiza zjawisk anomalnych}

W trakcie eksperymentów szczególna uwaga zostanie poświęcona trzem zjawiskom przejściowym, które mają istotny wpływ na zachowanie systemów \textit{SaaS} w środowisku chmurowym:

\begin{enumerate}
  \item \textbf{HPA Lag} (opóźnienie reakcji autoskalera) — opóźnienie $\sim$60--90 sekund między przekroczeniem progu metryki a faktycznym utworzeniem nowych Podów. Wynika ono z sumy: czasu scrapowania metryk przez \textit{Prometheusa} (15\,s) + czasu odświeżenia HPA (15\,s) + czasu tworzenia Podów i osiągnięcia gotowości. W tym okresie istniejące Pody są przeciążone, co może prowadzić do degradacji wydajności.
  
  \item \textbf{Cold Start Latency} (opóźnienie zimnego startu Podów) — nowo utworzone Pody potrzebują 15--30 sekund na osiągnięcie gotowości (pobranie obrazu z rejestru, uruchomienie kontenera, zaliczenie \textit{readiness probe}). W tym oknie czasowym ruch jest już dystrybuowany do nowych Podów przez \textit{LoadBalancer}, co powoduje podwyższone opóźnienia percentyla P99 lub błędy HTTP 502/503.
  
  \item \textbf{CPU Throttling Cascade} (kaskadowe błędy przy \textit{CPU throttlingu}) — gdy Pody osiągają limit CPU podczas gwałtownego wzrostu ruchu, mechanizm \textit{CPU throttling} w jądrze Linux (\textit{CFS bandwidth control}) ogranicza czas procesora. Powoduje to kolejkowanie żądań $\rightarrow$ wzrost czasu odpowiedzi $\rightarrow$ \textit{timeouty} w \textit{k6} (\texttt{dial: i/o timeout}) $\rightarrow$ błędy HTTP 502. Zjawisko to ma charakter kaskadowy, ponieważ HPA nie zdążył jeszcze dodać nowych replik (efekt \textit{HPA Lag}), tworząc ,,spiralę śmierci'', która ustępuje dopiero po osiągnięciu gotowości przez nowe Pody.
\end{enumerate}

\section{Podsumowanie architektury}

Zaprojektowany system badawczy umożliwia kontrolowane, powtarzalne eksperymenty porównujące skuteczność trzech strategii HPA dla dwóch profili obciążeniowych w środowisku GCP/GKE. Kluczowe decyzje architektoniczne:

\begin{enumerate}
  \item \textbf{Izolacja zmiennych} — każdy scenariusz testuje dokładnie jedną kombinację (profil $\times$ strategia), przy wyizolowanym generatorze obciążenia (osobna VM \textit{Compute Engine}) i ustabilizowanym klastrze (cooldown).
  \item \textbf{Powtarzalność} — każdy scenariusz jest powtarzany 5-krotnie, co umożliwia analizę statystyczną i eliminację wpływu anomalii.
  \item \textbf{Automatyzacja} — cały proces (przełączanie HPA, reset Deploymentów, uruchamianie k6, zbieranie wyników, eksport CSV z Grafany) jest zautomatyzowany przez skrypt \texttt{run-tests.sh}.
  \item \textbf{Odtwarzalność} — wszystkie manifesty \textit{Kubernetes}, skrypty k6 i konfiguracje są przechowywane w repozytorium, co umożliwia odtworzenie środowiska na dowolnym klastrze GKE.
  \item \textbf{Realizm wdrożeniowy} — wykorzystanie GCP LoadBalancer zamiast \texttt{port-forward} oraz dedykowanej maszyny \textit{Compute Engine} dla \textit{k6} odzwierciedla rzeczywiste warunki produkcyjne.
\end{enumerate}
